\chapter{Topology and Holomorphic maps}

\begin{lemma}[Normal form of holomorphic maps]
    Let \(f : X \to Y\) a non-constant holomorphic map. % between compact Riemann surfaces. ?
    Fix \(P \in X\). Then there exists \(m_P \in \mathbb{Z}_{\ge 1} \) such that
    for all charts \({(U_{2}, \varphi_{2} )} \ni f{(P)}\), \(\varphi {(f{(P)})} = 0\).
    Then one can find a chart \({(U_{1}, \varphi_{1})} \ni P\) with
    \(\varphi_{1}{(P)} = 0\) such that near 0, we have 
    \[
      \varphi_{2} \circ f \circ \varphi_{1}^{-1} {(z)} = z^{m_p}
    \]
    
\end{lemma}

\begin{proof}{}
    We start with two charts \({(U_{1}, \varphi_{1})} \ni P\) and \({(U_{2},
    \varphi_{2})} \ni f{(P)}\) such that \(\varphi_{1}{(P)} = 0\) and
    \(\varphi_{2}{(f{(P)})} = 0\).

    Near 0, we have that \(\varphi_{2} \circ f \circ \varphi_{1}^{-1}{(w)} =
    w^{m} S{(w)}\), with \(S {(0)} \neq 0\) and \(S\) analytic. We can use a
    complex logarithm to write \(S{(w)} = {\left( h{(w)} \right)} ^{m}\), hence
    \[
      \varphi_{2} \circ f \circ \varphi_{1}^{-1}{(w)} = {\left( \underbrace{w h{(w)}
      }_{=:\psi {(w)}} \right)} ^{m}
    \]
    but \(\varphi'{(0)} = h{(0)} \neq 0\) thus by local inverse theorem
    \(\varphi \) has a holomorphic inverse near 0, hence
    \[
        \varphi_{2} \circ f \circ {(\psi  \circ\varphi_{1})}^{-1}{(z)} = z^{m}
    \]
    and the chart is actually \(\psi \circ \varphi_{1}\) on a subset of \(U_{1}\) 
\end{proof}

\begin{definition}{}
    \(m_P := \mathrm{mult}_P{(f)}\) the \textbf{multiplicity}. If \(m > 1\),
    then \(P\) is called a \textbf{ramification point}.
\end{definition}

\begin{proposition}{}
    Ramification points corresponds to the set of vanishing points of \({\left(
    \varphi_{2} \circ f \circ \varphi_1^{-1}\right)}' \) at 0, hence it is in
    particular discrete.
\end{proposition}

\begin{theorem}{Definition of degree}
    Let \(f : X \to Y\), X, Y connected and \(f\) a \emph{proper} nonconstant
    holomorphic map.

    Then, for all \(y \in Y\),
    \[
      d_y{(f)} = \sum_{f{(x)} = y} \mathrm{mult}_x{(f)} \text{ is constant and
      finite}
    \]
    It is denoted \(\mathrm{deg}{(f)}\) and called the \textbf{degree} of \(f\).
\end{theorem}

\begin{proof}{}
    \(f\) being proper implies that
    \(f^{-1}{(\{y\} )}\) is compact. If it had infinitely many elements then I
    can construct a converging sequence. The local behavior on the accumulation
    point instead says that there should be only finitely many points in a
    neighborhood. Hence \(f^{-1}{(\{y\} )}\) is finite and thus also
    \(d_y{(f)}\).

    We want to show that \(y \mapsto d_y\) is locally constant. Because \(X\) is
    connected, this would imply \(y \mapsto d_y\) is constant. With this in
    mind, fix \(y \in Y\) and pick a chart \({(V, \psi)}\) at \(y\). Let
    \(f^{-1}{(\{y\} )} = \{x_{1}, \dots, x_{n}\} \). For all \(i = 1\dots n\)
    there exists \({(U_{i}, \varphi_{i})} \ni x_{i}\) such that \(\psi \circ f
    \circ \varphi_{i}^{-1}{(z)} = z^{m_{i}}\) and \(m_{i} =
    \mathrm{mult}_{x_{i}} {(f)}\). \\
    Now, for all \(y' \in V\),
    \[
      \sum_{x \in \bigcup_{i}^{n} U_{i} : f{(x)} = y'} \mathrm{mult}_x {(f)} = d_y{(f)} =
      \sum_{i=1}^{n} m_{i} 
    \]
    To show that \(d_{y'} {(f)} = d_y{(f)}\), we need to check that \[\{x \in X
        : f{(x)} = y'\} = \{x \in \bigcup_{i} ^{n} : f{(x)} = y'\} \]. We show
        this by contradiction. Suppose indeed there exists some \({(x'_k)}_{k
        \in \mathbb{N}} \in 
        f^{-1}{(\{y'_k\} )} \sminus \bigcup_{i} ^{n} U_{i} \) for all \(y'_k \to
        y\) as \(k\to \infty\) . By properness one
        can extract \(x'_k\) convergent.



\end{proof}

\begin{eser}{(T2 E1)}
    Soit \(n\ge 1\). Montrer que la courbe projective \(\mathcal{C} \subseteq
    \mathbb{P}^2{(\mathbb{C})} \) de Fermat d'équation homogène \(x^{n} + y^{n}
    + z^{n} = 0\) est non singulière.

    \tcblower

    Since the equation is symmetric in the coordinates it suffices to check one
    affine chart, say \(z \neq 0\). Hence \(g{(x, y)} = x^{n} + y^{n} + 1 = 0\) in
    \(\mathbb{C}^2\).
    \[
      \frac{\partial g}{\partial x} = n x^{n-1} = 0 \implies x = 0 \quad ; \quad
      \frac{\partial g}{\partial y} = n y^{n-1} = 0 \implies y = 0
    \]
    but \({(0, 0)} \not\in  \mathcal{C}\) 


    Another possibility is using Euler's formula, and is left as an exercise.
\end{eser}


\begin{eser}{(T2 E2)}
    Consider \(\mathcal{C} = \{y ^2 = f{(x)}\} \subseteq \mathbb{C}^2 \) where
    \(f \in  \mathbb{C}[x]\) of degree at least 3. Montrer les suivants:
\begin{enumerate}[label = (\arabic*)]
    \item affine curve \(\mathcal{C}\) is nonsingular \(\iff\) \(f\) doesn't
        have multiple roots
    \item  \(\mathcal{C} \subseteq \mathbb{P}^2{(\mathbb{C})} \) is a projective
        curve with affine equation \(y^2 = f{(x)}\) 
    \item The point \(\infty\) is nonsingular \(\iff\) \(d= 3\) 
\end{enumerate}
    \tcblower
\begin{enumerate}[label = (\arabic*)]
    \item \[
      \begin{cases}{}
          \frac{\partial }{\partial x} : f'{(x)} = 0\\
          \frac{\partial }{\partial y} : 2y = 0 \implies y = 0
      \end{cases}
    \]
    if \({(x, y)} \in \mathcal{C}\) then \(f{(x)} = 0\). Hence we get that
    \(f{(x)} = f'{(x)} = 0\) but then this means that \({(x, y)}\) is a singular
    point \(\iff\) \(y = 0\) and \(x\) is a multiple root of \(f\).
    \item Homogenizing we consider \(x = \frac{x_{1}}{x_{0}}\) and \(y =
        \frac{x_{2}}{x_{0}}\) and we get that \({\left(\frac{x_{2}}{x_{0}}\right)}^2 =
        f{\left(\frac{x_{1}}{x_{0}}\right)} = \frac{1}{x_{0}^{d}} F{(x_{1}, x_{0})}\) where
        \(F\) is the homogenization of \(f\).

        Now the equation becomes \(x_{0}^{d-2}x_{2} = F{(x_{1}, x_{0})}\).

        To now find the points at infinity consider \(x_{0}=0\) thus \(F{(x_{1},
        0)} = 0\) but \(F{(x_{1}, x_{0})} = a_d x_{1}^{d} + a_{d-1} x_{1}^{d-1} x_{0} + \dots 
        + a_{0} x_{0}^{d}\) hence \(x_{1}=0\). We get that the only point at
        infinity is \([0 : 0 : 1] \in  \mathcal{C}\) 
    \item Let's look at the affine chart \(\{x_{2} \neq 0\} \). Let's call \(u
        := \frac{x_{1}}{x_{2}}\) and \(v := \frac{x_{0}}{x_{2}}\).

        The equation in the coordinates \({(u,v)}\) is \(v^{d-2} = F{(u,v)}\).
        Let's check singularity:
        \[
          \begin{cases}{}
              \frac{\partial }{\partial u} : \frac{\partial F}{\partial u} = 0\\
              \frac{\partial }{\partial v} : {(d-2)}v^{d-3} - \frac{\partial
              F}{\partial v} = 0
          \end{cases}
        \]
        Since \(F\) is homogenous, then its partial derivates are also, hence
        \(\frac{\partial F}{\partial u} {(0, 0)} = \frac{\partial F}{\partial
        v}{(0,0)} = 0\). Then, if \(d > 3\) \({(0,0)}\) is a singular point and
        it is not if \(d = 3\). Note that in these coordinates \({(0,0)}\) is
        \(\infty\).

\end{enumerate}
\end{eser}

\begin{eser}{(T2 E3)}
    Consider \(a_{0}, \dots, a_{n} \ge 1\) integers. Consider the
    \emph{projective space à poids} the quotient \(X := \) 
    \[
      \mathbb{P}{(a_{0}, \dots, a_{n})} = {(\mathbb{C}^{n+1} \sminus \{0\} )} /
      \mathbb{C}^{\times } \text{ via the action } t{(x_{0}, .., x_{n})} = {(t
      ^{a_{0}} x_{0}, \dots, t ^{a_{n}} x_{n})}
    \]

    In the following consider \(a_{0} = a_{1} = \dots = a_{n}\) 

\begin{enumerate}[label = (\arabic*)]
    \item \(U_{0} = \{x_{0} \neq 0 \} \subseteq X \) is an affine chart and
        \(U_{0} \cong \mathbb{C}^{n}\) if \(a_{i} = 1\) for all \(i \ge  1\).
\end{enumerate}

\end{eser}

\begin{eser}{(T2 E5)}
    Let \(w_{1}, w_{2}\) a base of \(\mathbb{C}\) as a real vector space. Let
    \(f : \mathbb{C} \to \mathbb{C}\) be holomorphic such that \(f{(z + w_{i})}
    = \alpha_{i} f{(z)}\) for all \(z \in \mathbb{C}\) and \(i = 1, 2\). 

    Prove that \(f{(z)} = \lambda e^{\mu z}\) for some \(\lambda, \mu \in \mathbb{C}\) 

    \tcblower

    First note that
    \[
      f{(z + a_{1}\omega_{1} + a_{2}\omega_{2})} = \alpha_{1}^{a_{1}}
      \alpha_{2}^{a_{2}} f{(z)} \quad a_{1}, a_{2} \in \mathbb{Z}
    \]

    we may suppose that \(\alpha_{1}, \alpha_{2} \neq 0\) since otherwise \(f =
    0\). Note then that
    \[
      f'{(z + \omega_{i})} = \alpha_{i}f'{(z)}
    \]
    Let's then look at \(F = f'/f\) (knowing that the answer we are trying to find
    is the exponential) which is invariant under \(\Lambda\) and meromorphic. If
    \(f\) is invertible, the \(F\) is an holomorphic elliptic function, hence is
    constant. But then clearly \(f\) is of the form \(\lambda e ^{\mu z}\) for
    some \(\lambda and \mu\).

    We still need to show that \(f \neq 0\). If \(f{(z_{0})} = 0\). 
\end{eser}

\begin{corollary}{}
    Assume \(f : X \to Y\) holomorphic between \(X, Y\) compact connected
    Riemann Surfaces. Suppose \(f\) is non constant. Then 
    \[
      \mathrm{deg}{(f)} = 1 \iff f \text{ is biholomorphic}
    \]
\end{corollary}
\begin{proof}{}
    If the \(\mathrm{deg}{(f)} = 1\) then it must be injective by definition of
    degree.
\end{proof}

\begin{example}{}
    If \(X\) is compact connected and \(f : X \to \mathbb{C}\) is
    meromorphic with a unique simple pole. Then \(X \cong
    \widehat{\mathbb{C}}\).

    Indeed \(f^{-1}{(\{\infty\} )} = \{P\} \) hence \(\mathrm{deg}{(f)} = 1\) thus
    \(f\) is actually biholomorphic.
\end{example}

\section{\(\Delta\)-Complexes}
It doesn't work that well, that is why generally singular homology is used.

\begin{definition}{}
    A \textbf{simplex} is \(\Delta^{n} = \{{(x_{0}, \dots, x_{n})} \in \mathbb{R}^{n+1} :
    \sum_{i=0}^{n} x_{i} = 1, \, x_{i} \ge 0 \} \). \textbf{faces} are 
    \(F_{i} = \{{(\mathbf{x}  \in \Delta^{n} : x_{i} = 0)}\} \). There is a
    \textbf{face map} \(d_{i} : \Delta^{n-1} \to \Delta^{n}\) mapping \({(x_{0},
    \dots, x_{n-1} )}\) to \({(x_{0}, \dots, x_{i-1}, 0, x_{i}, \dots, x_{n-1}
    )}\). The \textbf{boundary map} is \(\partial \Delta^{n} = \bigcup_{i=0}
    ^{n}F_{i}\) 
\end{definition}

\begin{definition}{\(\Delta\)-complex}
    Let \({(\Delta, X)} = \{\sigma_{\alpha} : \Delta^{n_{\alpha} } \to
    X\}_{\alpha \in I}  \) with \(I\) \emph{finite}. These properties must be
    valid:
\begin{enumerate}[label = \arabic*.]
\item each map \(\sigma_{\alpha}|_{\mathring{\Delta}^{n_{\alpha} }}\) is injective and \(X = \bigcup_{\alpha
\in I} \sigma_{\alpha} {\left( \mathring{\Delta}^{n_{\alpha} } \right)} \) 
    \item  \(\forall \alpha \in I\), and \(\forall i\), \(\sigma_{\alpha} \circ
        d_{i} : \Delta^{n_{\alpha} - 1} \to X\) belong to \({(\Delta, X)}\) 
    \item The topology on \(X\) is such that \(A \subseteq X \) is open if and
        only if \(\forall \alpha \in I\), \(\sigma_{\alpha} ^{-1}{(A)} \subseteq
        \Delta^{n_{\alpha} }\) is open 
\end{enumerate}
    Hence each \(\sigma_{\alpha} : \overset{o}{\Delta^{n_{\alpha} }} \to
    \sigma_{\alpha} {\left( \mathring{\Delta}^{n_{\alpha} } \right)} \) is a
    homeo, and \(X \cong \coprod_{\alpha \in I} \Delta^{n_{\alpha} } /_\sim \)
    identification of faces.

\end{definition}
\begin{remark}{}
    the simplexes in the complex are not all the same.
\end{remark}

\begin{example}{}
    Let \(\tilde{X} = \Delta^2 \times \{1,2\} \) and \(\pi_{1} : \tilde{X} \to
    \Delta^2\) given by the projection \({(x, i)} \mapsto x\). Define next the
    relation
    \[
      x  \sim  y \iff \exists {(i, t)} \in \{0,1,2\} \times \Delta^{1} \text{ s.t. }
      \pi_{1}{(x)} = \pi_{1}{(y)} = d_{i}{(t)}
    \]
    Finally \(X = \tilde{X} /_\sim \) and let's note \([x,i] \in X\) the
    equivalence class.

    At this point there are the following
\begin{itemize}[label = --]
    \item[0-Simplexes:] \(\sigma_{i}^{0} : e_{i} \mapsto [e_{i}, 1]\) for \(i =
        0,1,2\) 
    \item[1-Simplexes:] \(\sigma_{i}^{1} : \Delta^{1} \ni t \mapsto [d_{i}{(t)},
        1\) for \(i = 0,1,2\) 
    \item[2-Simplexes:] \(\sigma_{i}^2 : \Delta^2 \ni t \mapsto [t, i] \) for
        \(i = 1,2\) 
\end{itemize}
\end{example}

\subsection{Homology of \(\Delta\)-simplexes}

Consider the complex \({(\Delta, X)} = \{\sigma_{\alpha} \} _{\alpha \in I} \)
and \(n := \max_{\alpha \in I} n_{\alpha} < \infty\). For all \(k = 0, \dots,
n\) let \(C_k\) be the free abelian group generated by \(\{\sigma_{\alpha}
{(\Delta^{n_{\alpha} })} : n_{\alpha} = k\} \).

Then the boundary map \(\partial_k : C_k \to C_{k-1} \) is \textbf{defined} by
\[
  \partial_k \{\sigma_{\alpha} {(\Delta^{n_{\alpha} })}\} = \sum_{i=0}^{k}
  {(-1)}^{i} \sigma_{\alpha} \circ F_{i}{\left( \Delta^{k-1} \right)}  
\]

\begin{example}{}
    Let \(\partial^{0} = \{1\} \), \(\partial^{1} = \mathrm{conv}{(e_{0},
    e_{1})} = [e_{0}, e_{1}]\) and \(\Delta^2 = \mathrm{conv} \{e_{0}, e_{1},
e_{2}\} = [e_{0}, e_{1}, e_{2}]\).

    Hence, the boundery maps give \(\partial_{1}\Delta^{1} = \{e_{1}\} -
    \{e_{0}\} \) and \(\partial_{2}\Delta^2 = [e_{1}, e_{2}] - [e_{0}, e_{2}] +
    [e_{0}, e_{1}]\) 
\end{example}

\begin{proposition}{}
    If we look at the sequence
    \[
      \{0\} \to C_{n} \overset{\partial_{n}}{\to} C_{n-1}
      \overset{\partial_{n-1} }{\to} \dots \overset{\partial_{2}}{\to} C_{1}
      \overset{\partial_{1}}{\to} C_{0}
    \]
    is a complex of abelian groups, i.e. \(\partial_{k-1} \circ \partial_k = 0\)
    for all \(k\) 
\end{proposition}

\begin{definition}{}
    We define the \textbf{homology groups}
    \[
      H_{k} ^{\Delta} = \frac{\mathrm{Ker}\partial_k}{\mathrm{Im} \partial_{k+1} }
    \]
\end{definition}

\begin{theorem}{}
    \(\Delta\)-complex homology is isomorphic to \emph{singular} homology.
\end{theorem}

\begin{definition}{}
    The Euler characteristic
    \[
      \rchi_{{(\Delta, X)}} = \sum_{i=0}^{n} {(-1)}^{i} \mathrm{rk}_{\mathbb{Z}}
      {\left( H_{i}^{\Delta} \right)} = \sum_{i=0}^{n} {(-1)}^{i}
      \mathrm{rk}_{\mathbb{Z}} {(C_{i})} 
    \]
    is a topological invariant
\end{definition}

\begin{theorem}{}
\begin{enumerate}[label = \arabic*)]
    \item All connected compact Riemann Surfaces are \(\Delta\)-complexes with \(n=2\) 
    \item If \(\mathcal{F} \subseteq X \) is finite, one can choose \({(\Delta, X)}\) such
        that \(\mathcal{F} \subseteq \bigcup_{n_{\alpha} = 0} \sigma_{\alpha}
        {\left( \Delta^{0} \right)}  \) 
    \item \(\rchi_X = \sum_{i=0}^{2} {(-1)}^{i} \mathrm{rk}{(C_{i})} = 2 - 2g \)
        is a topological invariant, where \(g\) is called \emph{genus}.
\end{enumerate}
\end{theorem}

\begin{theorem}{Riemann-Hurwitz}\label{thm:RH}
    Let \(X, Y\) be compact connected Riemann surfaces and \(f : X \to Y\) holomorphic
    non-constant. Then
    \[
      2g{(X)} - 2 = \mathrm{deg}{(f)} \, {(2 g{(Y)} - 2)} + \sum_{P \in X}
      {\left( \mathrm{mult}_P{(f)} - 1 \right)} 
    \]
\end{theorem}
\begin{proof}{}
    We want to avoid somehow the ramification points, since without them it
    would be a topological covering and much easier. 

    Let's call \(R{(f)} := \{x \in X : \mathrm{mult}_x{(f)} > 1\} \) and \(B{(f)}
    := f{(R{(f)})}\).
    We know that \(f : X \sminus f^{-1}\{B{(f)}\} \to Y \sminus B{(f)}\) is a
    topological covering of degree \(\mathrm{deg}\,f\). Take a
    \(\Delta\)-complex on \(Y\) such that \(B{(f)} \subseteq {(\Delta, Y)}_0 \).
    Remark that \(\overbrace{\sigma_{\alpha} {\left( \Delta^2 \right)} }^{o}
    \subseteq {(\Delta, Y)}_2 \) are simply connected.

    Above each \(\overbrace{\sigma_{\alpha} {\left( \Delta^2 \right)} }^{o}\),
    \(f\) is a trivial covering. Denoting by \(f^{-1}_{j,\alpha} :
    \overbrace{\sigma_{\alpha} {\left( \Delta^2 \right)}}^{o} \to X\) all \(d =
    \mathrm{deg}{(f)}\) local inverses of \(f\).

    The collection \(\left\{f_{j,\alpha} ^{-1} \circ \sigma_{\alpha} {\left(
    \mathring{\Delta}^2 \right)} \right\}_{j=1. .d}  \) can be completed as a
    \(\Delta\)-complex on \(X\). In particular, one need to add 0-simplexes and
    1-simplexes. We note that \(f^{-1}\{B{(f)}\} \subseteq {(\Delta, X)}_0 \).

    By the covering property, \(\# {(D, X)}_2 = d \# {(\Delta, Y)}_2\) and
    similarly \(\# {(\Delta, X)}_1 = d \# {(\Delta, Y)}_1\).

    The part to count are \(\# {(\Delta, X)}_0\). Pick \(q \in  {(\Delta,
    Y)}_0\) and \[\# f^{-1}{(\{q\} )} = \sum_{f{(p)} = q} 1 = \sum_{f(p) = q}
    {\left( 1 - \mathrm{mult}_{p} {(f)} + \mathrm{mult}_p {(f)} \right)} = d +
\sum_{f{(p)} = q} {\left( 1 - \mathrm{mult}_p{(f)} \right)}   \]
    Hence 
    \[
      \# {(\Delta, X)}_0 = \sum_{q \in  {(\Delta, Y)}_0} \# f^{-1}{(\{q\} )} = d
      \# {(\Delta, Y)}_0 + \sum_{q \in {(\Delta, Y)}_0} \sum_{f{(p)} = q}
      {\left( 1 - \mathrm{mult}_p{(f)} \right)}
    \]
    but the last double sum is just \(\sum_{p \in X} {\left( 1 - \mathrm{mult}_p
    {(f)}\right)} \). Finally,
    \[
      \rchi {(X)} = \# {(\Delta, X)}_2 - \# {(\Delta, X)}_1 + \# {(\Delta, X)}_0
      = d \rchi{(Y)} + \sum_{P \in X} {\left( 1 - \mathrm{mult}_p {(f)} \right)} 
    \]
\end{proof}

\subsection{Applications of R.H.}

\begin{theorem}{Placker's formula}
    let \(\mathcal{C}\) be a (smooth) projective curve of degree \(d\).
    then \[g{(\mathcal{C})} = \frac{{(d-1)}{(d-2)}}{2}\] 
\end{theorem}
\begin{proof}[proof sketch]
    Consider \(\mathcal{C} = \{{(x, y, z)} \in  \mathbb{P}^2\mathbb{C} : F{(x,
    y, z)} = 0\} \) and the projection \(\pi : \mathcal{C} \to
    \mathbb{P}^{1}\mathbb{C} \cong \widehat{\mathbb{C}}\) given by \(\pi{([x : y
    : z])} = [x : z]\). Then \(\pi\) is a holomorphic map of degree \(d\).

    Ramification points are the points \(\{[x : y : z] \in \mathcal{C} :
    \frac{\partial F}{\partial y} = 0\} \). Then from R.H. (theorem~\ref{thm:RH})
    we get that 
    \begin{equation}\label{eq:help1}
      2g{(\mathcal{C})} - 2 = d{(-2)} + \sum_{x \in \mathcal{C}} {\left(
      \mathrm{mult}_x{(\pi)} -1 \right)} 
    \end{equation}
    Then, by Bézout (a geometric algebra theorem),
    \begin{equation}\label{eq:help2}
      \sum_{x \in  \mathcal{C}} {\left( \mathrm{mult}_x{(\pi)} -1 \right)} =
      \deg {(F)} \cdot \deg {\left( \frac{\partial F}{\partial y} \right)} = d
      {(d-1)}
    \end{equation}
    By equations~\eqref{eq:help1} and~\eqref{eq:help2} we get that
    \[
      g{(\mathcal{C})} = \frac{1}{2} {\left( 2 - 2d + d{(d-1)} \right)} =
      \frac{{(d-1)}{(d-2)}}{2}
    \]
\end{proof}

